

\usepackage[spanish]{babel}
%\usepackage[spanish,es-tabla]{babel}

%\usepackage[spanish]{}
\usepackage [T1]{fontenc}
\setcounter{tocdepth}{2}

\newtheorem{defi}{{\bf Definición}}[chapter]
\usepackage{pifont}  % simbolos del 2023
\usepackage{amsmath}
\usepackage{enumerate} 
\usepackage{xcolor}




%  	\usepackage[color]{showkeys}    %  visualiza etiquetas
%%%%%%%%%%%%%%%%%%%%%%%%% hiper  vinculos %%%%%%%%%%%%%%%%%%%%%%%%%%%%%%%%%%
% \renewcommand{\chapter}{\LARGE{\bf {Capítulo}}}
\usepackage{hyperref}
\hypersetup{colorlinks , citecolor=Violet, 	linkcolor=black, 	urlcolor=black}
%  \usepackage[dvipdfm]{hyperref}
%  \tableofcontents


%\hypertarget{bar}
%Go to \hyperlink{bar}{bar graphic}


\usepackage{background} 

%%%%%%%%%%%%%%%%%%%%%%%%%%%%%%%%%%%%%
% Rotar texto
\usepackage{multirow}
\usepackage{rotating}





%%%%%%%%%%%%%%%%%%%%%%%%%%%%%%%%     Mini pagina    %%%%%%%%%%%%%%%%%%%%%%%%%%%%%%%%%%%%%
%% \usepackage[papersize={160mm,225mm},tmargin=15mm,bmargin=15mm,lmargin=15mm,rmargin=15mm]{geometry}

%    \usepackage[left=3cm,top=2cm,right=3cm,bottom=8.5cm]{geometry}
%     
%%%%%%%%%%%%%%%%%%%%%%%%%%%%%%%%%%%%%    %%%%%%%%%%%%%%%%%%%%%%%%%%%%%%%%%%%%%



%  \usepackage{python}
%  marca de agua

%%%%%%%%%%%%%%%%%%%%%%%%%  Marca de Agua  %%%%%%%%%%%%%%%%%%%%%%%%%%%%
% \usepackage{background} 
% \backgroundsetup{contents=Borrador, opacity=0.45, scale=10, color=blue, angle=45}

%\includegraphics{logo1} para figuras o graficos

%  https://www.youtube.com/watch?v=a3kbr4UtLXk
%%%%%%%%%%%%%%%%%%%%%%%%%  Marca de Agua  %%%%%%%%%%%%%%%%%%%%%%%%%%%%

% Marca de Agua
%\usepackage{background}


%\usepackage{listings}

\usepackage{dsfont}   %  Simbolo de  Numeros en latex

\usepackage{picinpar}  % Para abrir ventanas en parrafos

% Graficas en latex

\usepackage{graphicx} 

\graphicspath{{Imagenes/}} 

%	\usepackage{lscape}   %   pagina horizontal, pero no funciono, giro la figura
\usepackage{pdflscape}  %   pagina horizontal, si funciona


\DeclareGraphicsRule{.emf}{bmp}{}{}
\DeclareGraphicsExtensions{.pdf,.png,.jpg}     %solo para PDFLaTeX

\usepackage{wrapfig}

\usepackage[skins]{tcolorbox}

%	\usepackage{hyperref}         %%%% hiper enlaces 

\usepackage{fancybox}

%\addto\captionsspanish{
%	\renewcommand{\partname}{Libro}
%	\renewcommand{\chaptername}{Jornada}     }

\renewcommand{\partname}{Libro}

\renewcommand{\chaptername}{Capítulo}





%  \renewcommand{\thefootnote}{\arabic{footnote}}  Nota pie de pagina
%  \renewcommand{\chaptername}{Capítulo}



\usepackage{subcaption}




\usepackage{python}

%%%%%%%%%%%%%%%%%%%%%%%%%%%%%%%

%\definecolor{deepblue}{rgb}{0,0,0.5}
%\definecolor{deepred}{rgb}{0.6,0,0}
%\definecolor{deepgreen}{rgb}{0,0.5,0}
%\tiny       	5 5	6 6	6 6
%\scriptsize	    7 7	8	8
%\footnotesize








\usepackage{colortbl}
\definecolor{vino}{rgb}{0.5, 0.0, 0.13}

\definecolor{azul}{rgb}{0.1, 0.1, 0.9} 
\definecolor{azul}{rgb}{0.1, 0.1, 0.9} 
\definecolor{verde1}{rgb}{0.55, 0.71, 0.0}

\definecolor{verde2}{rgb}{0.01, 0.75, 0.24}

\definecolor{azul2}{rgb}{0.40, 0.56, 0.76}
\definecolor{amarillo1}{rgb}{0.98, 1, 0.40}


\usepackage{listings}
\usepackage{comment}

\lstset{language=Python,
	basicstyle=\ttfamily\scriptsize\scriptsize,  
	%	otherkeywords={self},             % Add keywords here
	%	keywordstyle=\color{blue},
    captionpos =b, %  Establece la posición de la leyenda del cuadro de código, pie de código, b= debajo	                  % 
	% commentstyle=\color{miverde},	
	emph={0,1,2,3,4,5,6,7,8,9},          % Custom highlighting
	emphstyle=  \color{black},         % \color{red!85!black},    % Custom highlighting style
	keepspaces= true,
	showspaces= false, 
	otherkeywords={0,1,2,3,4,5,6,7,8,9},              
%	rulecolor=  \color{green!85!black},%\color{green!85!black},% \color{red!70!blue},    %      \color{gray},	
	stringstyle= \color{green!85!black},    %    \color{red},
	% title=\lstname\color{green},
	keywordstyle= \color{red!70!blue},	%  \bf\color{blue!75!black},    % 
	%	frame=tb,                         % Any extra options here
	identifierstyle=  \color{black},   %  \color{blue},
  %  stringstyle=  \color{red},   %  \color{green!85!black}, 
   % otherkeywords=   \color{blue},             % Add keywords here
  %	keywordstyle=\color{blue},
	showstringspaces=false,
	commentstyle=\color{red}	
	      }
   




\lstset{literate=
	{á}{{\'a}}1 {é}{{\'e}}1 {í}{{\'i}}1 {ó}{{\'o}}1 {ú}{{\'u}}1
	{Á}{{\'A}}1 {É}{{\'E}}1 {Í}{{\'I}}1 {Ó}{{\'O}}1 {Ú}{{\'U}}1
	{à}{{\`a}}1 {è}{{\`e}}1 {ì}{{\`i}}1 {ò}{{\`o}}1 {ù}{{\`u}}1
	{À}{{\`A}}1 {È}{{\'E}}1 {Ì}{{\`I}}1 {Ò}{{\`O}}1 {Ù}{{\`U}}1
	{ä}{{\"a}}1 {ë}{{\"e}}1 {ï}{{\"i}}1 {ö}{{\"o}}1 {ü}{{\"u}}1
	{Ä}{{\"A}}1 {Ë}{{\"E}}1 {Ï}{{\"I}}1 {Ö}{{\"O}}1 {Ü}{{\"U}}1
	{â}{{\^a}}1 {ê}{{\^e}}1 {î}{{\^i}}1 {ô}{{\^o}}1 {û}{{\^u}}1
	{Â}{{\^A}}1 {Ê}{{\^E}}1 {Î}{{\^I}}1 {Ô}{{\^O}}1 {Û}{{\^U}}1
	{ã}{{\~a}}1 {ẽ}{{\~e}}1 {ĩ}{{\~i}}1 {õ}{{\~o}}1 {ũ}{{\~u}}1
	{Ã}{{\~A}}1 {Ẽ}{{\~E}}1 {Ĩ}{{\~I}}1 {Õ}{{\~O}}1 {Ũ}{{\~U}}1
	{œ}{{\oe}}1 {Œ}{{\OE}}1 {æ}{{\ae}}1 {Æ}{{\AE}}1 {ß}{{\ss}}1
	{ű}{{\H{u}}}1 {Ű}{{\H{U}}}1 {ő}{{\H{o}}}1 {Ő}{{\H{O}}}1
	{ç}{{\c c}}1 {Ç}{{\c C}}1 {ø}{{\o}}1 {å}{{\r a}}1 {Å}{{\r A}}1
	{€}{{\euro}}1 {£}{{\pounds}}1 {«}{{\guillemotleft}}1
	{»}{{\guillemotright}}1 {ñ}{{\~n}}1 {Ñ}{{\~N}}1 {¿}{{?`}}1 {¡}{{!`}}1 
}

%%=================================================================================%%
%%                        Pagina de fondo   \BgThispage                                          %%
%%=====================   Pagina da fondo    +=====================================%%

%\usepackage{scrextend}  %  Tamano de letra

%La opción pages puede ser all (para todo el documento) o some, para algunas partes del documento



% \usepackage[pages=some]{background}  %  some, y  no   all. Por pagina   pages=all  205



%    \BgThispage   por pagina la marca de aguaa

\backgroundsetup{     % aqui configuracion 
	scale=1, %escala de la imagen, es recomendable que sea del mismo tamaño que el pdf
	color=black, %fondo a usar para transparencia
	opacity=  1,  %       0.2, %nivel de transparencia
	angle=0, %en caso de querer una rotación
	contents={%
%	\includegraphics[width=1.6\linewidth, height=1.45\textheight]{Bill}
	\includegraphics[width=1.6\linewidth, height=1.45\textheight]{Texto1}
%   { \bf\Huge{ {\color{red} {\emph{T r a d u c c i ó n:  L i n i s \hspace{0.2cm} G u e r r e r o} } }} } %	\includegraphics[width=\paperwidth,height=\paperheight]{Traducción}%{Bill.jpg} %nombre de la imagen a utilizar como fondo  \LARGE
	 }}%
%  {Bill}

%    \BgThispage 

 %  Si la opción en pages=all, entonces automáticamente el fondo de asignará a toda la página, en caso 
 %   de usar some, entonces es necesario utilizar la etiqueta \BgThispage , para indicar que esa página llevará  % fondo, pero las otras páginas no.

 
  
  %
 
 
 